\documentclass[11pt,a4paper]{article}

\usepackage[T1]{fontenc}
\usepackage[utf8]{inputenc}
\usepackage[a4paper,margin=2.3cm]{geometry}
\usepackage{amsmath,amssymb}
\usepackage{graphicx}
\usepackage{booktabs}
\usepackage{xcolor}
\usepackage{listings}
\usepackage[colorlinks=true,linkcolor=blue!60!black,urlcolor=blue!60!black,citecolor=blue!60!black]{hyperref}
\usepackage[capitalise]{cleveref}
\usepackage{tcolorbox}
\tcbuselibrary{breakable,skins}
% microtype removed: lmodern is not installed in this TeX Live.
\usepackage{caption}
\captionsetup{font=small,labelfont=bf}

\definecolor{codebg}{HTML}{F5F5F5}
\definecolor{codekw}{HTML}{0000AA}
\definecolor{codecmt}{HTML}{006400}
\definecolor{codestr}{HTML}{A52A2A}

\lstdefinestyle{shell}{
  language=bash,
  backgroundcolor=\color{codebg},
  basicstyle=\ttfamily\small,
  breaklines=true,
  frame=single,
  framerule=0.4pt,
  framesep=4pt,
  xleftmargin=6pt,
  xrightmargin=6pt,
  commentstyle=\color{codecmt}\itshape,
  keywordstyle=\color{codekw}\bfseries,
  stringstyle=\color{codestr},
  showstringspaces=false,
}
\lstdefinestyle{py}{
  language=Python,
  backgroundcolor=\color{codebg},
  basicstyle=\ttfamily\small,
  breaklines=true,
  frame=single,
  framerule=0.4pt,
  framesep=4pt,
  xleftmargin=6pt,
  xrightmargin=6pt,
  keywordstyle=\color{codekw}\bfseries,
  commentstyle=\color{codecmt}\itshape,
  stringstyle=\color{codestr},
  showstringspaces=false,
}

\newtcolorbox{tipbox}{
  colback=blue!4!white, colframe=blue!50!black,
  arc=2pt, left=4pt, right=4pt, top=3pt, bottom=3pt,
  title=Tip,
  fonttitle=\bfseries,
  breakable
}
\newtcolorbox{warnbox}{
  colback=red!5!white, colframe=red!50!black,
  arc=2pt, left=4pt, right=4pt, top=3pt, bottom=3pt,
  title=Watch out,
  fonttitle=\bfseries,
  breakable
}

\title{\textbf{Student tutorial: a CPU-friendly pipeline for top-quark
tagging with classical ML and deep learning}}
\author{BNBWU--CERN FYP collaboration\\
        \small for BS Physics final-year students}
\date{\today}

\begin{document}
\maketitle
\begin{abstract}
This is the practical companion to the research article. It walks
you, step by step, through the codebase, the physics behind it, and
exactly which commands to type to reproduce every figure. By the end
you will (a) understand \emph{what} the pipeline does and \emph{why},
(b) be able to run it end-to-end on a laptop, and (c) be able to
modify any of the four models for your own experiments.
\end{abstract}

\tableofcontents
\vspace{0.4cm}

% ----------------------------------------------------------------
\section{What is this project, in one paragraph?}\label{sec:what}

At the LHC, a top quark almost always decays to a $W$ boson and a
$b$-quark; the $W$ then decays to two more quarks. When the top is
produced with high transverse momentum, all three of those quarks end
up inside a single, large-radius jet. The job of a ``top tagger'' is
to look at that jet and decide whether it really came from a top, or
just from an ordinary QCD multi-parton splash. We compare four
different ways to do that on a public benchmark dataset:
\begin{enumerate}
    \item A \textbf{boosted decision tree} on hand-engineered jet
          features (the classical baseline).
    \item A \textbf{convolutional neural network} on
          $(\eta,\phi)$ ``jet images''.
    \item \textbf{ParticleNet}, a graph neural network that treats
          the jet as an unordered cloud of particles.
    \item The \textbf{Particle Transformer}, the modern
          self-attention architecture used in CMS/ATLAS.
\end{enumerate}

The pipeline trains them, evaluates them, makes calibration plots and
saves everything to \texttt{results/}. The research angle is in
\texttt{run\_study.py} and asks: \emph{how much data does each model
need, and are its scores trustworthy probabilities?}

% ----------------------------------------------------------------
\section{A bit of physics background}\label{sec:physics}

\subsection*{The Standard Model, in 60 seconds}
The Standard Model has six quarks, six leptons and four gauge bosons.
The top quark is the heaviest, $m_t\!\simeq\!172.76$\,GeV. It is so
heavy that it decays before it can hadronise: more than 99\,\% of the
time, $t\to W b$, and the $W$ then decays either leptonically
(to $\ell\nu$) or hadronically (to $q\bar q$). In this project we
work with the \emph{fully hadronic} case
\[
   t \;\to\; W b \;\to\; (q\bar q)\,b ,
\]
which gives three quarks inside one ``boosted'' jet at high
$p_T$. The signature is a three-prong substructure and a jet mass
$\sim\!m_t$. QCD jets are dominantly one-prong and lighter --- this
is the physical signal/background asymmetry we are exploiting.

\subsection*{Why ``jet substructure''?}
At hadron colliders one cluster final-state particles into ``jets''
using algorithms like anti-$k_T$. Once a single boosted jet contains
the whole top decay, the \emph{internal} structure of the jet
(its mass, its $N$-subjettiness ratios $\tau_{21}$, $\tau_{32}$, its
energy correlation functions $C_2$) becomes the discriminating
information. The BDT baseline of this project uses exactly those
hand-engineered features.

\subsection*{Why deep learning, then?}
Hand-engineered features compress all the constituent information
into 10--20 numbers. A deep network is allowed to look at all the
constituent four-momenta and \emph{learn} its own features. The
question of this project is how much extra discrimination that
freedom buys --- and at what cost in training data and calibration.

% ----------------------------------------------------------------
\section{Project layout}\label{sec:layout}

\begin{lstlisting}[style=shell,caption={Top-level directory tree.}]
FYP_particle_physics/
+- config.py                  all hyperparameters & paths
+- download_data.py           Zenodo fetcher
+- run_full_pipeline.py       headline run (all 4 models)
+- run_study.py               data-efficiency + calibration sweep
+- requirements.txt
+- src/
|  +- data_utils.py           HDF5 reader, feature engineering
|  +- jet_cnn.py              CNN on jet images
|  +- particle_net.py         EdgeConv graph network
|  +- particle_transformer.py self-attention model
|  +- trainer.py              unified train/val/predict + BDT
|  +- evaluation.py           metrics + all publication plots
|  +- study_utils.py          data-efficiency + calibration plots
+- notebooks/                 exploratory analysis
+- docs/                      this tutorial + article + thesis
+- data/                      created by download_data.py
+- results/                   created at runtime
\end{lstlisting}

\begin{tipbox}
If you have never read this code before, the best entry points are
\texttt{run\_full\_pipeline.py} (orchestrator), \texttt{config.py}
(everything you can tune), and the \texttt{compute\_*\_features}
functions in \texttt{src/data\_utils.py} (the physics).
\end{tipbox}

% ----------------------------------------------------------------
\section{Setting up the environment}\label{sec:setup}

\subsection{macOS / Linux / WSL}
\begin{lstlisting}[style=shell]
# 1) clone / cd into the project
cd FYP_particle_physics

# 2) create a fresh virtual environment
python3 -m venv particle_venv
source particle_venv/bin/activate

# 3) upgrade pip and install dependencies
pip install --upgrade pip
pip install -r requirements.txt
\end{lstlisting}

\begin{warnbox}
Do not place the venv on a OneDrive or iCloud folder. Cloud
synchronisation will mark large binary files (e.g.\
\texttt{libtorch\_cpu.dylib}) as ``online-only'' placeholders, and
the operating system will then fail to memory-map them, causing
PyTorch to silently break at import time. Use a local, non-synced
folder.
\end{warnbox}

\subsection{Sanity check}
\begin{lstlisting}[style=shell]
python3 -c "import torch, xgboost, h5py; print(torch.__version__)"
\end{lstlisting}

If this prints a torch version, you are good.

% ----------------------------------------------------------------
\section{Downloading the dataset}\label{sec:download}

\begin{lstlisting}[style=shell]
python3 download_data.py     # answer 'y' at the prompt
\end{lstlisting}

This grabs three HDF5 files (about $1.6$\,GB total) from Zenodo and
drops them into \texttt{data/}. The script is resumable: if your
connection drops, just re-run it. If automatic download fails (e.g.\
proxy), download manually from
\href{https://zenodo.org/records/2603256}{zenodo.org/records/2603256}
and place \texttt{train.h5}, \texttt{val.h5}, \texttt{test.h5} into
\texttt{data/}.

% ----------------------------------------------------------------
\section{Running the headline pipeline}\label{sec:run}

The full headline run is a single command:
\begin{lstlisting}[style=shell]
python3 run_full_pipeline.py
\end{lstlisting}

That will:
\begin{enumerate}
    \item Read $5\!\times\!10^4$ jets per split from the HDF5 files.
    \item Compute particle-level and jet-level features, plus
          jet images.
    \item Train BDT, CNN, ParticleNet, Particle Transformer one
          after the other.
    \item Save checkpoints to \texttt{results/models/}.
    \item Generate all publication plots in
          \texttt{results/plots/} and a comparison table at
          \texttt{results/model\_comparison.json}.
\end{enumerate}

Useful flags:
\begin{lstlisting}[style=shell]
python3 run_full_pipeline.py --model bdt          # only BDT
python3 run_full_pipeline.py --model cnn          # only CNN
python3 run_full_pipeline.py --model particlenet  # only ParticleNet
python3 run_full_pipeline.py --model transformer  # only ParT
python3 run_full_pipeline.py --eval-only          # rebuild plots
                                                  # from saved ckpts
python3 run_full_pipeline.py --max-particles 50   # faster
\end{lstlisting}

\begin{tipbox}
The first time you run the pipeline, allow up to one hour on a
laptop CPU. Subsequent runs in \texttt{--eval-only} mode regenerate
plots from saved checkpoints in $\sim\!2$ minutes.
\end{tipbox}

\subsection{Stopping and resuming}\label{sec:resume}
The pipeline is fully checkpointed. You can Ctrl-C at any moment ---
between epochs, in the middle of a batch, even by closing the
laptop --- and pick up exactly where you stopped by re-running the
same command:
\begin{lstlisting}[style=shell]
python3 run_full_pipeline.py        # starts (or resumes)
# ... Ctrl-C ...
python3 run_full_pipeline.py        # resumes mid-CNN, then
                                    # continues with ParticleNet
                                    # and Particle Transformer
\end{lstlisting}

What is persisted as you go:
\begin{itemize}
  \item \texttt{results/models/<model>\_last.pt}: full training state
        (model, optimiser, scheduler, AMP scaler, epoch counter,
        history) snapshot after every epoch.
  \item \texttt{results/models/<model>\_best.pt}: best-val-AUC
        checkpoint, used for the final test prediction and the plots.
  \item \texttt{results/models/<model>\_done.json}: a small marker
        file written when training for that model has \emph{finished}
        (either max-epochs or early-stop).  Models whose marker
        exists are skipped on the next run.
  \item \texttt{results/predictions/<model>\_test.npz}: cached test
        predictions, used to skip the prediction step too.
  \item \texttt{results/<model>\_history.json}: epoch-by-epoch loss
        and AUC, updated every epoch (so the training-history plot
        reflects partial progress as well).
\end{itemize}

\noindent To force a fresh re-train, pass \verb|--force| (or delete
\texttt{results/models/}).  The same scheme applies to
\texttt{run\_study.py}: each (model, $N_{\text{train}}$) cell is
resumable independently.

% ----------------------------------------------------------------
\section{Running the research-angle study}\label{sec:study}

Once the headline run finishes, the research contribution comes from
\begin{lstlisting}[style=shell]
python3 run_study.py
\end{lstlisting}

This re-trains each model at $N_{\text{train}} \in
\{5\,000, 10\,000, 25\,000, 50\,000\}$ jets, measures test AUC and
background rejection at fixed signal efficiency, and computes the
probability calibration of each model with and without a
single-parameter temperature rescaling.

Selected flags:
\begin{lstlisting}[style=shell]
# pick the sizes to sweep
python3 run_study.py --sizes 5000 25000 50000

# only some models
python3 run_study.py --models bdt cnn

# cap epochs per training (default 15)
python3 run_study.py --epochs-cap 10
\end{lstlisting}

% ----------------------------------------------------------------
\section{Reading the plots}\label{sec:plots}

All plots land in \texttt{results/plots/}, in both PDF and PNG. The
intent of each:
\begin{description}
    \item[\texttt{jet\_substructure.pdf}]\quad Signal/background
        distributions of all 13 engineered observables. The
        \emph{physics check} -- if this looks wrong, the dataset
        loading is broken.
    \item[\texttt{average\_jet\_images.pdf}]\quad The mean
        $p_T$-weighted image for top vs.\ QCD. You should
        \emph{see} the three-prong topology in the difference
        panel.
    \item[\texttt{roc\_comparison.pdf}]\quad ROC and background
        rejection vs.\ signal efficiency. The headline result.
    \item[\texttt{score\_distributions.pdf}]\quad Per-model
        classifier output distributions. Look for the
        characteristic ``peak at 0 / peak at 1'' shape that
        well-trained binary classifiers produce.
    \item[\texttt{confusion\_matrices.pdf}]\quad At threshold 0.5.
        Off-diagonal entries are misclassifications.
    \item[\texttt{training\_history.pdf}]\quad Train/val loss and
        AUC vs.\ epoch. Use it to diagnose overfitting (train
        loss falling while val loss climbing) and to set the
        early-stop patience.
    \item[\texttt{bdt\_importance.pdf}]\quad XGBoost gain-based
        feature importance. You should see jet mass, $\tau_{32}$
        and constituent count among the top three.
    \item[\texttt{attention\_maps.pdf} /
          \texttt{cls\_attention.pdf}]\quad Particle Transformer
        interpretability.
    \item[\texttt{learning\_curves.pdf}]\quad Output of
        \texttt{run\_study.py}: AUC and rejection vs.\ training-set
        size.
    \item[\texttt{calibration.pdf}]\quad Reliability diagrams and
        ECE before/after temperature rescaling.
\end{description}

% ----------------------------------------------------------------
\section{How to modify the pipeline}\label{sec:modify}

\paragraph{Try a new dataset subset size.}
Edit \texttt{config.py}: change \texttt{SUBSET\_TRAIN /
SUBSET\_VAL / SUBSET\_TEST}.

\paragraph{Try a deeper CNN.}
Edit \texttt{CNN\_CONFIG[``filters'']} in \texttt{config.py}, or
the residual structure in \texttt{src/jet\_cnn.py}.

\paragraph{Change the kNN of ParticleNet.}
\texttt{PARTICLENET\_CONFIG[``k\_neighbors'']} in \texttt{config.py}.

\paragraph{Add more pairwise features for the Transformer.}
The pairwise features are computed in
\texttt{compute\_pairwise\_features\_batch} inside
\texttt{src/data\_utils.py}. Add another channel there (e.g.\ the
invariant mass $m_{ij}$) and bump
\texttt{PARTFORMER\_CONFIG[``pair\_features'']} accordingly.

\paragraph{Add a new model.}
\begin{enumerate}
    \item Create a new file \texttt{src/my\_model.py} with a
          PyTorch class implementing \texttt{forward(features,
          mask)} (and \texttt{count\_parameters}).
    \item Register a hyperparameter dict in \texttt{config.py}.
    \item Add a new branch in \texttt{run\_full\_pipeline.py}
          following the \texttt{ParticleNet} branch as a template.
\end{enumerate}

% ----------------------------------------------------------------
\section{Troubleshooting}\label{sec:trouble}

\begin{description}
\item[\texttt{ImportError: dlopen ... libtorch\_cpu.dylib}]\quad
The venv is on a OneDrive/iCloud-synced path. Re-create it on a
local, non-synced folder (see \cref{sec:setup}).

\item[\texttt{zsh: segmentation fault python3 run\_full\_pipeline.py}]\quad
macOS OpenMP collision between NumPy's \texttt{libiomp5} and
XGBoost's \texttt{libomp}. \texttt{run\_full\_pipeline.py} already
sets \texttt{KMP\_DUPLICATE\_LIB\_OK=TRUE} and pins XGBoost to one
thread; if it persists, export
\texttt{OMP\_NUM\_THREADS=1} in your shell as well.

\item[\texttt{TypeError: a bytes-like object is required, not 'str'}
on \texttt{pd.read\_hdf}]\quad This pandas/PyTables bug is patched
in \texttt{src/data\_utils.py}: it reads HDF5 directly with
PyTables and decodes block-column metadata by hand.

\item[\texttt{FileNotFoundError: data/train.h5}]\quad Run
\texttt{python3 download\_data.py}.

\item[\texttt{Killed} / OOM during data prep]\quad Lower the
\texttt{SUBSET\_*} values in \texttt{config.py}, or pass
\texttt{--max-particles 50}.
\end{description}

% ----------------------------------------------------------------
\section{Suggested student exercises}\label{sec:exercises}

\begin{enumerate}
\item \textbf{Mass-decorrelation.}  Modify the BDT to mask out the
      \texttt{jet\_mass} feature. By how much does AUC drop? Plot
      the new ROC against the old one.
\item \textbf{Robustness to constituent count.}  Restrict
      \texttt{--max-particles} to $20$, $50$, $100$. How does the
      Transformer scale compared to ParticleNet?
\item \textbf{Try a fifth model.}  Implement a simple Deep Sets
      baseline (per-particle MLP + sum pooling) and add it to the
      learning-curve plot.
\item \textbf{Re-do calibration with isotonic regression.}
      Implement post-hoc isotonic calibration (using
      \texttt{sklearn.isotonic.IsotonicRegression}) and compare its
      ECE against the temperature rescaling.
\item \textbf{Pile-up robustness.}  Smear constituent $p_T$'s with
      a Gaussian kernel and re-evaluate. Which model degrades
      least?
\end{enumerate}

\section*{Further reading}
\begin{itemize}
    \item Kasieczka \emph{et al.}, ``The Machine Learning landscape
          of top taggers'', SciPost Phys.\ 7 (2019) 014 ---
          benchmark survey.
    \item Qu and Gouskos, ``ParticleNet: Jet Tagging via Particle
          Clouds'', Phys.\ Rev.\ D 101 (2020) 056019.
    \item Qu, Li and Qian, ``Particle Transformer for Jet
          Tagging'', ICML 2022, arXiv:2202.03772.
    \item Guo \emph{et al.}, ``On Calibration of Modern Neural
          Networks'', ICML 2017, arXiv:1706.04599.
\end{itemize}

\end{document}
